% ============================================================================
% Глава 7: Изгледи и Шаблонизиране (View & Templating)
% ============================================================================
\section{Изгледи и Шаблонизиране (View \& Templating)}
\label{sec:views-templating}

С развитието на екосистемата на Coroute възникна необходимост от по-директна техническа интеграция за генериране на HTML съдържание на сървъра (Server-Side Rendering). За тази цел е създаден модулът \texttt{view/}, предоставящ мощна, но лесна за използване система за шаблони.

\subsection{Интеграция с Inja}
Модулът използва \texttt{Inja} като вътрешен двигател за шаблоните. Inja е съвременна C++ библиотека, която имплементира синтаксиса на Jinja (популярен в Python екосистемата) и се интегрира безпроблемно с JSON структурата на \texttt{nlohmann::json}. Този избор гарантира гъвкавост, защото шаблоните могат лесно да итерират през масиви, да проверяват условни логики директно и да форматират данни.

\subsection{Компоненти на View Системата}

В основата на архитектурата стоят няколко ключови компонента, които заедно осигуряват предсказуем жизнен цикъл на шаблона:

\subsubsection{\texttt{ViewTemplates} кеширане}
Класът \texttt{ViewTemplates} действа като глобален регистър (или singleton асоцииран с основния \texttt{App} контекст). Неговата задача е да зареди съдържанието на шаблоните от файловата система или паметта по време на инициализация, да ги парсне чрез Inja и да съхрани компилираното представителство (AST - Abstract Syntax Tree) в кеша. 

\begin{lstlisting}[language=C++, basicstyle=\small\ttfamily, caption={Конфигурация на ViewTemplates}]
ViewTemplates::get().add_template_file(
    "auth_info", 
    "templates/auth_info.html"
);
\end{lstlisting}

Това предварително компилиране означава, че по време на обработка на HTTP заявка не се извършва дисков I/O, нито лексикален анализ на шаблона — само оценяване на АСТ дървото с подадения контекст. Това драстично повишава пропускателната способност.

\subsubsection{\texttt{ViewRouteInfo} за конфигурация}
\texttt{ViewRouteInfo} предоставя структура за дефиниране на метаданните на даден изглед -- например кое име на шаблон извиква и кои динамични параметри очаква от контролера (handler).

\subsubsection{Лесно конструиране чрез \texttt{ResponseBuilder}}
Интеграцията с мрежовия слой е опростена чрез разширения в обекта за отговори. Модулите могат да върнат готов HTML отговор, използвайки хелпър функции, които автоматично подават \texttt{nlohmann::json} контекст към кеширания шаблон:

\begin{lstlisting}[language=C++, basicstyle=\small\ttfamily, caption={Връщане на рендиран шаблон}]
app.get("/profile", [](Request& req) -> Task<Response> {
    nlohmann::json context;
    context["username"] = "admin";
    context["last_login"] = "2024-03-01";
    
    // Автоматично рендира auth_info шаблона и задава Content-Type: text/html
    co_return ResponseBuilder::render_view("auth_info", context);
});
\end{lstlisting}

\subsection{Изпълнение на шаблонизираните маршрути}

При получаване на заявката към \texttt{/profile}, handler-ът изгражда съответния JSON контекст. След това, извикването на \texttt{render\_view} подава заявката към Inja интерпретатора, който намира AST-то в паметта, асоциирано с ключа \texttt{auth\_info}. Интепретаторът генерира крайния HTML низ. \texttt{ResponseBuilder} автоматично конструира валиден HTTP 200 OK отговор с хедър \texttt{Content-Type: text/html; charset=utf-8} и го изпраща към клиента.

Този механизъм позволява разработването на класически SSR (Server-Side Rendered) уебсайтове и портали директно в Coroute, съвместно с JSON-базираните RESTful API endpoints, без да се налага въвеждането на допълнителни външни уеб сървъри.
